\documentclass[11pt]{article}

\setlength{\parindent}{0pt}
\usepackage{xltxtra}
\usepackage{hyperref}
\hypersetup{hidelinks}
\usepackage{url}
\urlstyle{tt}
\usepackage{xcolor}
\definecolor{CVBlue}{RGB}{23,110,191}
\usepackage{calc}
\usepackage{graphicx}
\usepackage{tikz}
\usetikzlibrary{calc}
\usepackage{fontspec}
\usepackage{xeCJK}
\usepackage{enumitem}
\CJKsetecglue{} %% 取消中文与数字之间的间隙

%% 主文档字体设置
\setmainfont[
    Path = fonts/Main/,
    Extension = .otf,
    BoldFont = texgyretermes-bold.otf, % 加粗字体
]{texgyretermes-regular.otf} % 正文字体

% 中文字体设置
\setCJKmainfont[
    Path = fonts/hansans/,
    Extension = .ttf,
    BoldFont = NotoSansSC-Bold.ttf, % 加粗字体
]{NotoSansSC-Regular.otf} % 正文字体

\usepackage{relsize}
\usepackage{xspace}
\usepackage{fontawesome} 

% A4纸，上下左右边距
\usepackage[
    a4paper,
    left=1.2cm,
    right=1.2cm,
    top=1.5cm,
    bottom=1cm,
    nohead
]{geometry}

\renewcommand{\baselinestretch}{1.5} % 行间距设为1.5

\usepackage{titlesec}
\setlist{noitemsep} % 取消列表项间的额外间距
\setlist[itemize]{topsep=0.25em, leftmargin=*}
\setlist[enumerate]{topsep=0.25em, leftmargin=*}

% --- 用于控制【不同项目之间】的垂直距离 ---
\newlength{\interProjectSpacing}
\setlength{\interProjectSpacing}{0.9em} 
\newcommand{\projectsep}{\vspace{\interProjectSpacing}}

% --- 用于控制【项目标题】与下方【项目描述】的距离 ---
\newlength{\intraProjectTitleSep}
\setlength{\intraProjectTitleSep}{0.4em} 
\newcommand{\titlebreak}{\\[\intraProjectTitleSep]}

% --- 用于控制【项目描述】与下方【要点列表】的距离 ---
\newlength{\intraProjectListTopSep}
\setlength{\intraProjectListTopSep}{0.2em} 

% =======================================================================

\titleformat{\section}         
{\large\bfseries\raggedright} 
{}{0em}                      
{}                           
[{\color{CVBlue}\titlerule}]  
\titlespacing*{\section}{0cm}{*1.6}{*1.2}

\begin{document}
\pagenumbering{gobble}

%%%% 如需照片可解除此段注释并调整图片名称
% \begin{tikzpicture}[remember picture, overlay] 
%     \node[anchor = north east] at ($(current page.north east)+(-2cm,-1.2cm)$) {\includegraphics[height=3cm]{avatar.jpg}};
% \end{tikzpicture}%

\centerline{\LARGE\bfseries{史阳春}} 
\vspace{0.5em}
\centerline{\normalsize{2004年3月 \quad 汉族 \quad 安徽省合肥市}} 
\vspace{0.3em}
\centerline{\normalsize{\faPhone\ 19826568639 \quad \faEnvelopeO\ \href{mailto:582323815@qq.com}{582323815@qq.com}}} 
    
\section{\makebox[\widthof{\faGraduationCap}][c]{\color{CVBlue}\faGraduationCap}\ 教育背景}    
\textbf{安徽工业大学} \hfill 2023.09 -- 2027.06\\[0.5em] 
工业工程 \quad 本科
\begin{itemize}[nosep]
    \item \textbf{成绩}：平均学分绩：83.35/100 \quad GPA：3.04/4
    \item \textbf{主修课程}：基础工业工程、人因工程、生产计划与控制、运筹学、生产系统建模与仿真、精益生产、六西格玛管理、物流与供应链管理
\end{itemize}

\section{\makebox[\widthof{\faUsers}][c]{\color{CVBlue}\faUsers}\ 项目经历}

\textbf{不确定加工时间的电机生产节能调度优化研究} \quad 负责人 \hfill 2025.04 -- 2027.04 \titlebreak
项目描述：研究通过改进智能算法构建的多目标调度模型，所构建的优化调度模型可直接应用于电机生产场景。通过优化的算法模型，可降低设备空载率及无效能耗，助力企业年均碳排放量削减。
\begin{itemize}[nosep, topsep=\intraProjectListTopSep]
    \item \textbf{项目成果}：1. 2025年中国工业工程与精益管理创新赛国家二等奖；2. Applied Sciences在投论文一篇。
\end{itemize}

\section{\makebox[\widthof{\faBriefcase}][c]{\color{CVBlue}\faBriefcase}\ 实习经历}

\textbf{大陆马牌轮胎（中国）有限公司} \quad IE实习生 \hfill 2025.06 -- 2025.09 \titlebreak
\begin{itemize}[nosep, topsep=\intraProjectListTopSep]
    \item \textbf{数据分析}：测算热准备及成型 workload，收集生产数据，使用 Excel/Power BI 等工具分析数据，生成相应的可视化报表，以支持改善决策。
    \item \textbf{质量管理}：参与质量体系管理，协助分析缺陷原因，推动防错装置 (EPD) 验证流程的标准化。
    \item \textbf{物流与仿真}：负责硫化车间生胎转运的物流优化，利用 MTM 对设备进行改善设计，仿真建模改善物流路径和设备运作工艺。
    \item \textbf{产线改善}：跟踪产线异常情况（不良品，设备故障），分析并协助资源进行改善优化。
    \item \textbf{路径规划}：协助设备供应商规划 AGV 运行路径与扫描识别范围，并绘制相应的 CAD。
\end{itemize}

\section{\makebox[\widthof{\faTrophy}][c]{\color{CVBlue}\faTrophy}\ 竞赛经历}
\begin{itemize}[nosep, parsep=0.5ex]
    \item \textbf{国家级三等奖} \quad 第十九届“东风日产杯”清华IE亮剑全国工业工程创新大赛 (队长) \hfill 2024.11
    \item \textbf{国家级三等奖} \quad 第二十届“东风日产杯”清华IE亮剑全国工业工程创新大赛 (队长) \hfill 2025.12
    \item \textbf{华东赛区三等奖} \quad 工业工程与精益管理创新赛 (队长) \hfill 2025.10
\end{itemize}

\section{\makebox[\widthof{\faCogs}][c]{\color{CVBlue}\faCogs}\ 应用技能}
\begin{itemize}[nosep, parsep=0.5ex]
    \item \textbf{外语能力：} 通过英语四级考试 (454)，具备良好的口语能力。
    \item \textbf{专业证书：} 获得 MTM 体系认证证书。
    \item \textbf{专业技能：} 熟悉仿真建模理论，熟练使用 Flexsim 软件；了解工程制图理论，熟练 CAD 操作。
    \item \textbf{编程与数据：} 掌握 Python 编程语言；了解数据库技术，熟悉 SQL 编程。
    \item \textbf{办公工具：} 熟练操作 Excel、PPT 等办公软件。
\end{itemize}

\end{document}
